% swift-property-wrappers.tex — @propertyWrapper under the hood: wrappedValue,
% projectedValue ($), init(wrappedValue:), the synthesized _name storage,
% composition, parameter/local wrappers, the enclosing-instance subscript
% (@Published), how @State / @Binding / @AppStorage are built, limits.
% SwiftUI usage is on swiftui-state — this folio is the mechanism.
% Sources (checked 2026-10-06):
%   swift-evolution SE-0258 (property wrappers, 5.1: wrappedValue, projectedValue, _name
%     storage, composition — the outermost wrapper owns $, the static enclosing-instance
%     subscript) · SE-0293 (wrappers on parameters + init(projectedValue:), 5.5) — checked
%     against the evolution index
%   docs.swift.org/swift-book Properties › Property Wrappers (local wrappers allowed)
%   developer.apple.com/documentation/combine/published ("publishes … in the property's
%     willSet block"; the publisher emits the current value on subscription) ·
%     /swiftui/binding · /swiftui/state/init(initialvalue:)
% @labels: area=swift kind=concept platform=apple topic=language,ui discussion=162
% @tags: propertywrapper, wrappedvalue, projectedvalue, published, state, binding, appstorage, dynamicproperty, enclosing-instance, objectwillchange
\documentclass[8pt]{extarticle}
\usepackage{printup-sheet}
\usepackage{array}

\lstdefinelanguage{SwiftSheet}{
  morekeywords={protocol,class,final,struct,enum,func,var,let,init,case,static,
    subscript,get,set,nonmutating,private,extension,
    if,else,return,guard,self,nil,some,switch,default,in,true,false},
  sensitive=true, morecomment=[l]{//}, morecomment=[s]{/*}{*/}, morestring=[b]",
  literate={??}{{\hbox{?}\hbox{?}}}2 {==}{{\hbox{=}\hbox{=}}}2
           {!=}{{\hbox{!}\hbox{=}}}2 {->}{{\hbox{-}\hbox{>}}}2}

\newcommand\arr{\hbox{-}\hbox{>}}

\tikzset{
  sb/.style={box, font=\ttfamily\scriptsize, inner sep=2pt, minimum height=5mm, align=left},
  src/.style={sb, draw=sheetBrown, fill=sheetBrown!8},
  gen/.style={sb, draw=sheetGreen, fill=sheetGreen!8},
  ext/.style={sb, draw=sheetOrange, fill=sheetOrange!10},
  lbl/.style={font=\tiny, text=black!75, inner sep=1pt, align=center},
  hd/.style={font=\bfseries\small, anchor=west},
}

\begin{document}

\sheettitle{Property wrappers — what the compiler writes for you}{swift · folio}

\oneliner{\texttt{@propertyWrapper} (SE-0258, Swift 5.1) is a type with a \texttt{wrappedValue};
annotating \texttt{@W var x} makes the compiler store a hidden \textbf{\texttt{\_x: W}} and turn
\texttt{x} into a \textbf{computed property} forwarding to \texttt{\_x.wrappedValue} —
plus \textbf{\texttt{\$x}} for \texttt{projectedValue}. SwiftUI's \texttt{@State},
\texttt{@Binding}, \texttt{@AppStorage} and Combine's \texttt{@Published} are just such types.}

\vspace{3pt}
\noindent\begin{tikzpicture}[sheet]
  \draw[sheetGrey!40] (6.9,1.5) -- (6.9,-2.5);
  \draw[sheetGrey!40] (11.55,1.5) -- (11.55,-2.5);
  % ── A: desugaring ──
  \node[hd] at (0,1.35) {\textcolor{sheetBlue}{A} you write $\to$ the compiler writes};
  \node[src, anchor=west] (s) at (0,0.75) {@Clamped(0...100) var volume = 50};
  \node[gen, anchor=west, text width=62mm] (g) at (0,-0.6) {private var \_volume =
    Clamped(wrappedValue: 50, 0...100)\\[1pt]
    var volume: Int \{\\
    \ \ get \{ \_volume.wrappedValue \}\\
    \ \ set \{ \_volume.wrappedValue = newValue \} \}\\[1pt]
    var \$volume: \ldots\ \{ \_volume.projectedValue \}};
  \draw[flow] (s.south) ++(-1.6,0) -- ++(0,-0.33);
  \node[lbl, anchor=north west, align=left, text width=64mm] at (0,-1.75) {\texttt{= 50} feeds
    \texttt{init(wrappedValue:)}, attribute args follow it. \texttt{\_x} is private; \texttt{\$x} exists
    only if \texttt{projectedValue} does. Memberwise init takes \texttt{volume: Int} when the wrapper
    has \texttt{init(wrappedValue:)}, else the wrapper type.};
  % ── B: @Published ──
  \node[hd] at (7.0,1.35) {\textcolor{sheetBlue}{B} \texttt{@Published} reaches its owner};
  \node[src, anchor=west] (w) at (7.0,0.8) {vm.name = "Ann"};
  \node[gen, anchor=west, text width=40mm] (sub) at (7.0,-0.1) {static subscript(\_enclosingInstance: vm,\\
    \ wrapped: \textbackslash.name, storage: \textbackslash.\_name)};
  \node[ext, anchor=west] (ow) at (7.0,-1.05) {vm.objectWillChange.send()};
  \node[ext, anchor=west] (pub) at (7.0,-1.65) {\$name subscribers get "Ann"};
  \draw[flow] (w.south) ++(-0.5,0) -- ++(0,-0.22);
  \draw[flow] (sub.south) ++(-1.4,0) -- ++(0,-0.3);
  \node[lbl, anchor=west, text=sheetRed] at (7.0,-2.2) {both in \textbf{willSet}: \texttt{vm.name} still old here\\
    no \texttt{self} in a struct $\to$ \texttt{@Published} is class-only};
  % ── C: @State ──
  \node[hd] at (11.65,1.35) {\textcolor{sheetBlue}{C} \texttt{@State}: handle vs storage};
  \node[src, anchor=west, text width=19mm, fill=white, opacity=0.6, text opacity=0] at (11.83,0.83) {Counter()\\\_n: State<Int>};
  \node[src, anchor=west, text width=19mm, fill=white, opacity=0.8, text opacity=0] at (11.74,0.74) {Counter()\\\_n: State<Int>};
  \node[src, anchor=west, text width=19mm, fill=white] (v1) at (11.65,0.65) {Counter()\\\_n: State<Int>};
  \node[ext, anchor=west, text width=16mm] (gr) at (14.55,0.55) {graph storage\\n = 3};
  \draw[flow] ([xshift=1mm]v1.east) -- node[lbl, below]{finds} (gr.west |- v1.east);
  \node[gen, anchor=west] (bd) at (11.65,-0.65) {\$n: Binding<Int>};
  \draw[flow] (gr.south) |- node[lbl, right, pos=0.3]{get/set} (bd.east);
  \node[lbl, anchor=west] at (11.65,-1.45) {view struct is rebuilt every update;\\
    storage lives in SwiftUI, keyed by identity.\\\texttt{nonmutating set}: writable from \texttt{body}.\\
    Initial value used only the \textbf{first} time.};
\end{tikzpicture}

\begin{multicols}{2}
\raggedright

\section{How it works}
\begin{itemize}
  \item \textbf{Required:} \texttt{var wrappedValue}. Optional: \texttt{projectedValue} (any
        type — \texttt{Binding}, a publisher, the wrapper itself),
        \texttt{init(wrappedValue:)} (enables \texttt{= default}),
        \texttt{init(projectedValue:)} (enables passing \texttt{\$arg}).
  \item \textbf{Composition} \texttt{@A @B var x: T}: storage \texttt{\_x: A<B<T>>}, access
        \texttt{\_x.wrappedValue.wrappedValue}; outer to inner, left to right;
        \texttt{\$x} is the \textbf{outermost} wrapper's projection.
  \item \textbf{Where:} stored \texttt{var} properties of types; \textbf{locals} (Swift 5.4);
        function and closure \textbf{parameters} (SE-0293, 5.5) — that is how
        \texttt{ForEach(\$items) \{ \$item in TextField("", text: \$item.name) \}} works
        (\texttt{Binding} has \texttt{init(projectedValue:)}).
  \item \textbf{Enclosing-instance subscript} (underscored, used by \texttt{@Published}):
        if the wrapper declares \texttt{static subscript(\_enclosingInstance:wrapped:storage:)},
        the compiler routes access through it with the \textbf{owning object} and two key
        paths — so the wrapper can reach \texttt{self} (e.g. \texttt{objectWillChange}).
        Class-only: \texttt{Published.wrappedValue} is marked unavailable,
        \emph{``@Published is only available on properties of classes''}.
  \item \textbf{\texttt{@State}}: a \texttt{DynamicProperty} holding a handle; SwiftUI finds it
        on the view, allocates storage in its graph (per view \emph{identity}), updates the
        handle before \texttt{body}. \texttt{\$n} = \texttt{Binding}.
  \item \textbf{\texttt{@Binding}}: a struct of \texttt{get}/\texttt{set} closures to someone
        else's storage; \texttt{@dynamicMemberLookup}, so \texttt{\$user.name} is a
        \texttt{Binding<String>}. No \texttt{init(wrappedValue:)} $\to$ the child's
        memberwise init takes a \texttt{Binding}: \texttt{Child(n: \$count)}.
  \item \textbf{\texttt{@AppStorage("key")}}: a \texttt{DynamicProperty} over
        \texttt{UserDefaults}; the view re-renders when that key changes; \texttt{\$} =
        \texttt{Binding}. Bool/Int/Double/String/URL/Data + \texttt{RawRepresentable}.
\end{itemize}

\section{Limits}
\begin{itemize}
  \item Not on \texttt{let}, \texttt{lazy}, \texttt{weak}/\texttt{unowned}, computed
        properties, or a property with its own \texttt{get}/\texttt{set}
        (\texttt{willSet}/\texttt{didSet} are allowed).
  \item A protocol cannot require a wrapper — only the plain property.
  \item Inline storage (\texttt{@Clamped}): the setter is \texttt{mutating} — needs a
        \texttt{var} owner; external storage (\texttt{@State}) uses \texttt{nonmutating set}.
\end{itemize}

\columnbreak

\section{Example — a wrapper + the enclosing subscript}
\begin{lstlisting}[language=SwiftSheet]
@propertyWrapper struct Clamped<V: Comparable> {
  private var value: V; let range: ClosedRange<V>
  init(wrappedValue v: V, _ r: ClosedRange<V>) {
    range = r; value = Self.clamp(v, r) }       // = 50 lands here
  var wrappedValue: V {
    get { value }
    set { value = Self.clamp(newValue, range) } } // mutating set
  var projectedValue: ClosedRange<V> { range }   // $volume
  static func clamp(_ v: V, _ r: ClosedRange<V>) -> V {
    min(max(v, r.lowerBound), r.upperBound) }
}
// the shape @Published uses to reach its owner:
static subscript<Owner: AnyObject>(_enclosingInstance o: Owner,
  wrapped: ReferenceWritableKeyPath<Owner, Value>,
  storage: ReferenceWritableKeyPath<Owner, Self>) -> Value
\end{lstlisting}

\section{Traps}
\begin{itemize}
  \trap{\texttt{self.count = 5} in a View's \texttt{init} does not seed \texttt{@State} —
        write \texttt{\_count = State(initialValue: 5)}; and only the first init wins.}
  \trap{\texttt{@State var items = load()} runs \texttt{load()} on \emph{every} view init,
        result discarded after the first.}
  \trap{\texttt{\$vm.name.sink \{ \}} reading \texttt{vm.name} inside gets the \textbf{old}
        value (willSet); use the closure's parameter.}
  \trap{Two wrappers: \texttt{\$x} projects only the outermost one.}
  \trap{Wrappers hide semantics (persistence, threading) at the call site — great for
        cross-cutting concerns, costly when overused.}
\end{itemize}

\section{Remember}
\textbf{``\texttt{x} is a computed façade, \texttt{\_x} is the wrapper, \texttt{\$x} is the
projection. Storage inside = value; storage outside = survives the struct.''}


\end{multicols}

\noindent{\footnotesize\color{sheetGrey}\textit{Related:} swiftui-state (which wrapper to use) ·
observation-and-combine (\texttt{@Published}, \texttt{@Observable} macro) · macros-and-result-builders ·
swiftui-performance-identity · swift-initialization (memberwise)}

\end{document}
